\section{Optimal transport of the sampling}\label{sec:ot}

\subsection{Equidistribution and its indeterminacy}

Fix a reference torus $\Xi=\T^d$ carrying a fixed mesh (uniform, or a fixed hierarchy). A discrete sheet is
parameterised over $\Xi$ through a \emph{label map} $T:\Xi\to\Q$. Let $M:\Q\to(0,\infty)$ be a \emph{monitor
density}, $\int_\Q M\,\dd q=1$, expressing where resolution is wanted (\cref{sec:monitor}). The
\emph{equidistribution principle} asks that reference cells of equal volume carry equal monitor mass:
\begin{equation}\label{eq:equi}
T\pf\lambda=\nu_M:=M\,\dd q\qquad\Longleftrightarrow\qquad M(T(\xi))\,\det DT(\xi)=1 .
\end{equation}
For $d\ge2$ this is one scalar equation for $d$ unknown functions: if $T$ solves \eqref{eq:equi} so does
$s\circ T$ for every $\nu_M$-preserving diffeomorphism $s$. The indeterminacy is exactly the fibre $\Diff_{\nu_M}$
of \cref{rem:noether}; its elements shear and rotate cells without changing their monitor mass, and they are
harmful (they degrade mesh quality and enlarge the constants of every interpolation estimate) but invisible to \eqref{eq:equi}.

\subsection{The transport problem and its well-posedness}

\begin{definition}[sampling transport problem]\label{def:ot}
Given $M$, find
\begin{equation}\label{eq:monge}
T=\arg\min\Big\{\int_\Xi d_{\T^d}(\xi,T(\xi))^2\,\dd\xi\ :\ T\pf\lambda=\nu_M\Big\},
\end{equation}
where $d_{\T^d}$ is the flat distance on the torus (we identify $\Xi$ and $\Q$ as flat tori).
\end{definition}

\begin{theorem}[existence, uniqueness, structure, regularity]\label{thm:ot}\status{established
\cite{Brenier1991,Cordero1999,McCann2001,Caffarelli1992}}
Let $M\in C^{k,\alpha}(\Q)$, $k\ge0$, $0<\alpha<1$, with $0<m_-\le M\le m_+$. Then \eqref{eq:monge} has a unique
minimiser. It has the form
\begin{equation}\label{eq:brenier}
T(\xi)=\xi+\nabla\varphi(\xi)\pmod{\Z^d},\qquad \varphi\ \text{periodic},\qquad \tfrac12\abs{\xi}^2+\varphi\ \text{convex},
\end{equation}
$\varphi\in C^{k+2,\alpha}$ solves the Monge--Amp\`ere equation
\begin{equation}\label{eq:ma}
M\big(\xi+\nabla\varphi(\xi)\big)\,\det\big(I+D^2\varphi(\xi)\big)=1 ,
\end{equation}
and $T$ is a $C^{k+1,\alpha}$ diffeomorphism isotopic to the identity. Conversely a periodic $\varphi$ with
$I+D^2\varphi>0$ solving \eqref{eq:ma} gives the minimiser.
\end{theorem}

The selection made by \eqref{eq:monge} has three properties that matter here. \emph{(a) Canonical gauge.} Among all
equidistributing maps the optimal one is the unique gradient map (Brenier's polar factorisation: every
equidistributing map is $T\circ s$ with $s$ measure preserving); its displacement is curl free, which removes the
harmful fibre directions. \emph{(b) No tangling.} $I+D^2\varphi>0$ makes $T$ a diffeomorphism; in computations the
smallest eigenvalue of $I+D^2\varphi$ is a certificate. \emph{(c) Memorylessness.} $T$ depends only on the current
monitor, not on the history of the mesh; composing increments would not have this property (a composition of
gradient maps is not a gradient map) and would accumulate shear. Optimal-transport mesh adaptation with these
properties is well established for PDE meshes in physical space
\cite{BuddWilliams2006,BuddHuangRussell2009,Delzanno2008,Weller2016}; the present use differs only in \emph{where}
it acts (label space of a Lagrangian sheet) and in what it is allowed to touch (\cref{prop:sampling}).

\begin{proposition}[$d=1$: monotone rearrangement]\label{prop:ot1d}\status{established}
For $d=1$ the equidistributing orientation-preserving maps of the circle form a one-parameter family
$T_\tau=C^{-1}(\cdot+\tau)$, $C(q)=\int_{0}^{q}M$, and \eqref{eq:monge} selects the offset $\tau$ by a
one-dimensional minimisation \cite{DelonSalomonSobolevski2010}. Optimal transport therefore adds nothing to
equidistribution in $d=1$ beyond fixing a rotation; the prototype pins the label $q=0$ instead.
\end{proposition}

\paragraph{Linearisation and solution.} With $F(\varphi)=\log\det(I+D^2\varphi)+\log M(\xi+\nabla\varphi)$,
\begin{equation}\label{eq:malin}
F'(\varphi)[\psi]=\tr\big((I+D^2\varphi)^{-1}D^2\psi\big)+\nabla\log M(T)\cdot\nabla\psi ,
\end{equation}
a uniformly elliptic operator on admissible $\varphi$, equal to the Laplacian at $\varphi=0$, $M=1$; for
$M=1+\varepsilon m$ one has $\Delta\varphi=-\varepsilon m+O(\varepsilon^2)$. The prototype solves \eqref{eq:ma} by a
damped Newton iteration with \eqref{eq:malin} inverted by GMRES preconditioned with the FFT inverse Laplacian and
a backtracking line search that maintains $I+D^2\varphi>0$ (\texttt{solve\_ot} in \texttt{prototype/sheet2d.py}).
Measured behaviour is reported in \cref{sec:num2d}.

\subsection{Sampling is not mass}

\begin{proposition}[the physical measure is a passenger]\label{prop:sampling}\status{new derivation; elementary}
Let $T_0,T_1:\Xi\to\Q$ be two label maps and $F$ a sheet with mass measure $\mu$.
\begin{enumerate}[label=(\roman*),leftmargin=2.2em,itemsep=1pt]
\item If the reference density is updated as a pull-back, $\tilde\mu_i=T_i^*\mu$, both discretisations represent the
same measure $F\pf\mu$, for arbitrary monitors. The reference density is $\tilde\rho=(\rho_0\circ T)\det DT=(\rho_0/M)\circ T$:
cell masses necessarily vary by the factor $\max(M/\rho_0)/\min(M/\rho_0)$.
\item If instead the masses stay attached to the moved sample points (the reference density is kept,
$\tilde\mu_1:=\tilde\mu_0$), the represented measure becomes $F\pf(R\pf\mu)$ with $R=T_1\circ T_0^{-1}$, and
\begin{equation}\label{eq:w2}
W_2\big(F\pf R\pf\mu,\,F\pf\mu\big)\le\Lip(F)\,\norm{R-\id}_{L^2(\mu)},
\end{equation}
with equality to leading order for small displacements along a single stream. This is an $O(1)$ modification of the
distribution function unless $R\pf\mu=\mu$.
\item Equal-mass sampling and adaptive sampling are incompatible: $\tilde\rho$ constant forces $T\pf\lambda\propto\mu$,
leaving only the gauge freedom $\Diff_\mu$.
\end{enumerate}
\end{proposition}
\begin{proof}
(i) is \cref{thm:gauge}(i). (ii) $(F\circ T_1)\pf(T_0^{-1})\pf\mu=F\pf R\pf\mu$; the coupling $(F\circ R,F)\pf\mu$ gives
the bound. (iii) is \eqref{eq:equi} with $M=\rho_0/\mu(\Q)$.
\end{proof}

This is the precise sense in which the optimal transport of \cref{def:ot} must be distinguished from the uses of
optimal transport that \emph{do} move mass in cosmology: Monge--Amp\`ere--Kantorovich reconstruction and its
semi-discrete descendants \cite{FrischEtAl2002,BrenierEtAl2003,LevyEtAl2021} transport the mass density in
\emph{configuration space} between initial and final times and are restricted to single-stream (convex-potential)
maps. Here the transported measure is Lebesgue measure on a reference torus, the target is a sampling density on
label space, and the phase-space measure is invariant by construction, in any number of streams.

\subsection{A symplectically invariant monitor}\label{sec:monitor}

What should be equidistributed? Phase space has no canonical metric ($x$ and $p$ have different units), so
curvature-based monitors depend on an arbitrary scaling. The action supplies an invariant alternative.

\begin{lemma}[edge action defect]\label{lem:edge}\status{new derivation; verified symbolically}
Let $(F,\Sigma)$ be an exact $C^4$ sheet, $q_0\in\Q$, $e\in\R^d$ a label-space edge vector of length $h$, and
$z_0=F(q_0)$, $z_1=F(q_0+e)$. The \emph{edge action defect}
\begin{equation}\label{eq:edgedefect}
\eps_e:=\Sigma(q_0+e)-\Sigma(q_0)-\tfrac12(p_0+p_1)\cdot(x_1-x_0)
\end{equation}
is the signed symplectic area between the sheet arc over the edge and its chord, and
\begin{equation}\label{eq:edgeexp}
\eps_e=-\tfrac1{12}\,C(e,e,e)+O(h^4),
\end{equation}
with $C$ the cubic form \eqref{eq:cubic} at $q_0$.
\end{lemma}
\begin{proof}
$\Sigma(q_0+e)-\Sigma(q_0)=\int_0^1p\cdot\partial_sx\,\dd s$ along $s\mapsto F(q_0+se)$, and the trapezoid term is the
integral of $\theta$ along the chord, so $\eps_e=\oint\theta$ over the closed curve arc--minus--chord, which is the
signed area by Stokes. Expanding $x,p$ to third order in $s$,
$\tfrac12(p_0+p_1)\cdot(x_1-x_0)-\int_0^1p\cdot x'\,\dd s=\tfrac1{12}(x'\cdot p''-p'\cdot x'')+O(h^4)$ with
derivatives along $e$ at $s=0$, which is $\tfrac1{12}\om(\partial_eF,\partial_e^2F)$.
\end{proof}

$\eps_e$ is computable from vertex data alone, exactly (no finite differences), and is invariant under affine
symplectic maps of phase space (it is an area); it is the local quantity a piecewise-linear
representation gets wrong. Requiring $\abs{\eps_e}$ to be equidistributed for edges of local length $h(q)$, with
node density $h^{-d}$ and $N$ nodes, the standard variational argument (minimise $\int\norm{C}h^{3-d}\dd q$ subject
to $\int h^{-d}\dd q=N$) gives
\begin{equation}\label{eq:monitor}
M\ \propto\ \norm{C}^{\,d/(3+d)}\quad(L^1\text{-optimal}),\qquad M\ \propto\ \norm{C}^{\,d/3}\quad(\text{max-optimal}).
\end{equation}

\paragraph{The case $d=1$: equi-affine sampling and a complexity functional.} Here $C=\kappa\,\dd q^3$ and the
max-optimal monitor is $M\propto\abs{\kappa}^{1/3}$, i.e.\ \emph{uniform sampling in equi-affine arclength}, the
arclength that is invariant under area-preserving affine maps of the phase plane. This connects to the classical
theory of best polygonal approximation of convex curves in the area metric \cite{McClureVitale1975}. Writing
\begin{equation}\label{eq:complexity}
\begin{gathered}
\mathcal A(t):=\int_\Q\abs{\kappa(q,t)}^{1/3}\dd q,\qquad
\eps_{\rm equi}=\frac{\mathcal A^3}{12N^3},\qquad
\eps_{\rm unif}=\frac{\abs{\Q}^3\max\abs\kappa}{12N^3},\\
G(t):=\frac{\eps_{\rm unif}}{\eps_{\rm equi}}=\frac{\abs{\Q}^3\max_q\abs\kappa}{\mathcal A(t)^3}\ \ge1 ,
\end{gathered}
\end{equation}
$\mathcal A$ is a symplectic invariant measuring the intrinsic cost of representing the sheet by chords, and $G$ is
the gain of adaptive over fixed (uniform-in-label) sampling at equal $N$ ($G\ge1$ by H\"older). Adaptivity cannot
reduce $\mathcal A$; it can only remove $G$. The measured growth of $\mathcal A$ and $G$ is reported in
\cref{tab:complexity}: in the collapse and chaotic tests $\mathcal A$ grows roughly linearly while $G$ grows
exponentially.

\begin{remark}[regularisation and regularity of the monitor]\label{rem:monreg}
\eqref{eq:monitor} degenerates where $C=0$ and ignores higher-order error terms; practical monitors add a floor and
arclength/turning-angle terms (prototype: \texttt{monitor\_density}). More importantly, \emph{the remap error
inherits the regularity of $T$}: an interpolation estimate of order $r$ in the reference coordinates involves
$\norm{F\circ T}_{W^{r,\infty}}$, hence $r$ derivatives of $T$, hence $r-1$ derivatives of $M$ by \cref{thm:ot}. In
the $d=2$ prototype a monitor with a kink ($\norm{C}$ at a zero of $C$) smoothed only at the mesh scale reduced
the observed remap order to about one; a monitor mollified at a fixed label-space scale restored the design
order (\cref{sec:num2d}). In $d=1$ the prototype interpolates directly on the non-uniform label nodes, so the
regularity of $T$ does not enter. Finally, optimal transport controls a scalar density only: it cannot align cells
with the anisotropy of the tensor $C$; anisotropic (metric-based) refinement as in ColDICE addresses a complementary need.
\end{remark}
